\section{Secondary derivations}\label{app:derivations}

The main proofs are in the body: walk counting (\cref{thm:walkcount}), the walk-to-path lemma (\cref{lem:walkpath}), and the connectivity test (\cref{prop:pair}, \cref{prop:iplusa}, \cref{thm:test}).
This appendix proves the two statements about distance that the body uses, and describes the finite checks.

\subsection{Distance: shortest walks, trails, and paths}\label{app:distance}

\plainwords The distance between two houses was defined with the shortest trip that never drives a road twice, in the textbook's word, or the shortest route, in ours.
We show that the shortest of every kind of trip has the same length.
So it does not matter which word we use.
Then we check the three facts about distance that the textbook lists.

\begin{proposition}\label{prop:shortest}
For vertices $u,v$ of a simple graph the length of a shortest $(u,v)$-walk, of a shortest $(u,v)$-trail, and of a shortest $(u,v)$-path are equal.
Each is $+\infty$ when the corresponding set is empty, and the three sets are empty together.
Hence $d(u,v)$ in \cref{def:distance} is the same whether ``path'' there is read as trail or as path.
\end{proposition}

\begin{proof}
Every path is a trail (\cref{prop:pathtrail}) and every trail is a walk, so the minimum over walks is at most the minimum over trails, which is at most the minimum over paths.
By \cref{lem:walkpath} a shortest walk contains a path of at most the same length, so the minimum over paths is at most the minimum over walks.
Therefore the three minima are equal.
If there is no walk, there is no trail and no path; if there is a walk, \cref{lem:walkpath} gives a path.
So the three sets are empty together.
\end{proof}

Grinberg makes the same remark: a shortest walk from $u$ to $v$, if it exists, is also a shortest path \cite{grinberg2025graph}.

\begin{proposition}[properties of distance]\label{prop:distprops}
The three properties listed after \cref{def:distance} follow from the definition.
\end{proposition}

\begin{proof}
By \cref{prop:shortest} we may compute $d$ with shortest walks.
\emph{Property 1.}
Lengths are $0,1,2,\dots$ or $+\infty$, so $d\ge 0$.
The walk consisting of $u$ alone has length $0$, so $d(u,u)=0$.
If $d(u,v)=0$ there is a walk of length $0$ from $u$ to $v$; it is a single vertex, so $v=u$.
\emph{Property 2.}
If $u_0,\dots,u_k$ is a $(u,v)$-walk, then $u_k,\dots,u_0$ is a $(v,u)$-walk of the same length, since each step is an edge and edges are unordered.
So the sets of walk lengths from $u$ to $v$ and from $v$ to $u$ are equal, and so are their minima.
\emph{Property 3.}
If $d(u,w)$ or $d(w,v)$ is $+\infty$, the right side is $+\infty$ and the inequality holds.
Otherwise take a shortest $(u,w)$-walk and a shortest $(w,v)$-walk and join them at $w$: the result is a $(u,v)$-walk of length $d(u,w)+d(w,v)$, and a shortest $(u,v)$-walk is no longer.
\end{proof}

\subsection{Finite checks by computer}\label{app:checks}

The following checks were run by the scripts of \cref{app:repro}.
They are tests of statements already proved, not proofs.
\begin{itemize}
\item \emph{Walk counting.} For every labelled simple graph with $N=2,3,4,5$ vertices ($2$, $8$, $64$, $1024$ graphs), all walks were enumerated by brute force and compared with the entries of $A^k$ ($k=1,\dots,5$ for $N\le 4$ and $k=1,\dots,3$ for $N=5$): $40$, $360$, $5120$, and $76800$ entries respectively, all equal.
\item \emph{Connectivity test.} For every labelled simple graph with $N=2,\dots,6$ vertices ($2$, $8$, $64$, $1024$, $32768$ graphs, of which $1$, $4$, $38$, $728$, $26704$ are connected) the sum test, the $(I+A)^{N-1}$ test, the component statement of \cref{cor:components}, and the repeated-squaring version of \cref{prop:iplusa} all agree with an independent breadth-first search (\cref{tab:exhaustive}).
\item \emph{Walks, trails, paths.} The script for \cref{sec:walkdefs} classifies every example of that section by applying the definitions literally, and tests the construction of \cref{lem:walkpath}, \cref{prop:pathtrail}, \cref{cor:bound}, and \cref{prop:shortest} on all short walks of $B$, $T$, $H$ and on random graphs; the saved output records no failure.
\end{itemize}
